\section{Average Baby PIH}
\label{sec:avg}

\noindent Let us recall the average Baby PIH conjecture. 
\begin{hypothesis}[Average Baby PIH]
    Given a 2CSP instance $\Pi=(X,\Sigma,\Phi)$, we say a multi-assignment $\sigma$ $r$-average-list satisfies $\Pi$ if $\sum_{x \in X}|\sigma(x)|\le r\cdot |X|$, and for every constraint $\phi_i$, there exists $u \in \sigma(w_{i,1})$ and $v \in \sigma(w_{i,2})$, such that $(u,v) \in R_i$.

    Average Baby PIH states that for any constant $r>1$ and any computable function $f$, no algorithm can given as input a 2CSP instance $\Pi=(X,\Sigma,\Phi)$, distinguish between the following two cases in $f(|X|)\cdot |\Sigma|^{O(1)}$ time:
    \begin{itemize}
        \item (Completeness) $\Pi$ is satisfiable.
        \item (Soundness) $\Pi$ is not $r$-average-list satisfiable.
    \end{itemize}
\end{hypothesis}

\subsection{A Counter Example for Direct Product Construction}

We use the following counter example to show that the Direct Product construction does not give us Average Baby PIH. Specifically, for any $t>0$ and $\varepsilon>0$, there exists an 2CSP instance which is not satisfiable but its $t$-wise direct product is $(1+\varepsilon)$-average-list satisfiable.

\begin{example}\label{ex:1}
    The 2CSP instance $\Pi=(X,\Sigma,\Phi)$ is defined as follows.
    \begin{itemize}
        \item $X=\{x_1,\ldots,x_n\}$.
        \item $\Sigma_{x_1}=\{2,\ldots,n\}$, and for every $i \in \{2,\ldots,n\}$, $\Sigma_{x_i}=\{1\}$.
        \item For every $i \in \{2,\ldots,n\}$, there is a constraint $\phi_i=\langle w_i,R_i\rangle \in \Phi$, where
        \begin{itemize}
            \item $w_i=(x_1,x_i)$.
            \item $R_i=\{(j,1) | j \ne i\}$.
        \end{itemize}
    \end{itemize}
    \end{example}
    $\Pi$ is not satisfiable since no value for $x_1$ can satisfy all constraints. Specifically, for $i \in \{2,\ldots,n\}$, $x_1=i$ would violate constraint $\phi_i$.

    However, for every integer $t>0$, $\Pi^{\odot t}$ can be list satisfied by the following multi-assignment $\sigma$. For every $S \in \binom{X}{t}$, 
    \begin{itemize}
        \item if $x_1 \notin S$, define $\sigma(S)$ to be the single assignment that maps every $x_i \in S$ to 1;
        \item if $x_1 \in S$, then for every $2 \le j \le 2t$ with $x_j \notin S$, let $\sigma(S)$ contain an assignment that maps $x_1$ to $j$, and maps everything else to 1.
    \end{itemize}
    
    It's easy to see $\sigma$ satisfies all constraints induced by $S$. It remains to show that $\sigma$ is consistent on every different $S,T \in \binom{X}{t}$. 
    
    If $x_1 \notin S$ or $x_1 \notin T$, $\sigma(S)$ and $\sigma(T)$ are trivially consistent since every variable in $S \cap T$ is always assigned 1. Now suppose $x_1 \in S\cap T$, we have $|S \cup T| \le 2t-1$. By Pigeonhole Principle, there is a variable $x_j$ with $j \le 2t$, which is neither in $S$ nor in $T$. Thus the assignment that maps $x_1$ to $j$ and everything else to $1$ appears in both $\sigma(S)$ and $\sigma(T)$, proving that $\sigma$ list-satisfies $\Pi^{\odot t}$. The total size of the lists is \[
    \begin{aligned}
        \frac{1}{\binom{|X|}{t}}\sum_{x \in X} |\sigma(x)| & \le \frac{\binom{|X|-1}{t}}{\binom{|X|}{t}}\cdot 1+\frac{\binom{|X|-1}{t-1}}{\binom{|X|}{t}}\cdot 2t \\ 
        & = \left(1-\frac{t}{|X|}\right)\cdot 1+\frac{t}{|X|}\cdot 2t \\
        & \le 1+\frac{2t^2}{|X|},
    \end{aligned}\]
    which is smaller than $1+\varepsilon$ when $|X|$ goes to infinity.
    

\subsection{Towards the Inapproximability of \kexact}
In this subsection, we show that a slightly strengthened version of Average Baby PIH, namely, Average Baby PIH with rectangular relations, implies constant inapproximability of \kexact{} problem. The latter, to our best knowledge, is currently not known under $\wone \neq \fpt$. Formally, we have the following theorem:
\begin{theorem}\label{thm:exactcover}
    Suppose Average Baby PIH holds even when the 2CSP instance has rectangular relations (see Definition \ref{def:rect}), then for any constant $r \ge 1$ and any function $f$, no algorithm can approximate \kexact{} problem within factor $r$ in $f(k)\cdot n^{O(1)}$ time. Specifically, no algorithm can given a \ksetcover{} instance $\Pi=(\mathcal S,U)$ with size $n$, distinguish between the following two cases in $f(k)\cdot n^{O(1)}$ time:
    \begin{itemize}
        \item There exists $k$ sets in $\mathcal S$ which form a partition of $U$.
        \item The union of any $r\cdot k$ sets in $\mathcal S$ is not $U$.
    \end{itemize}
\end{theorem}
\begin{proof}
We reduce a 2CSP instance $\Pi=(X,\Sigma,\Phi)$ with rectangular relations to a \ksetcover{} instance $\Pi'=(\mathcal S,U)$ with $k=|X|$ in the following way. For every $(x,v) \in X \times \Sigma$, we build a set $S_{x,v}$. For each constraint $\phi_i=\langle (w_{i,1},w_{i,2}),R_i\rangle$ in $\Phi$, let $\Gamma_i$ be the underlying set in the rectangular relation $R_i$ and let $\pi_i,\sigma_i:\Sigma\to \Gamma_i$ be the underlying mappings. We build a $(r \cdot k,|\Gamma_i|)$-set gadget $(\mathcal M^{(i)},C_1^{(i)},\ldots,C_m^{(i)})$. For every $(a,b) \in R_i$, we add the set $C_{\pi_i(a)}^{(i)}$ to $S_{w_{i,1},a}$, and add the set $\overline{C_{\sigma_i(b)}^{(i)}}$ to $S_{w_{i,2},b}$. Let the final universe $U$ be the union of every $\mathcal M^{(i)}$.

%We build a $(2,|C_i|)$-hypercube partition system $\left(\mathcal M^{(i)}, \{P^{(i)}_{x,y}\}_{x \in C_i, y \in [2]}\right)$ and add the subset $P^{(i)}_{\pi_i(a),1}$ to $S_{w_{i,1},a}$, $P^{(i)}_{\sigma_i(b),2}$ to $S_{w_{i,2},b}$ for every $(a,b)\in R_i$. Let the final universe $U$ be the union of every $\mathcal M^{(i)}$. 

In the completeness case, let $\sigma:X \to \Sigma$ be a satisfying assignment of $\Pi$. It is easy to see the $k$ sets $\{S_{x,\sigma(x)}\}_{x \in X}$ cover each element of the universe exactly once.

In the soundness case, let $\mathcal S' \subseteq S$ be a collection of sets that covers $U$. Assuming $|\mathcal S'| \le r \cdot k$, we claim the multi-assignment $\sigma$, which maps $x \in X$ to $\{v \in \Sigma \mid S_{x,v} \in \mathcal S'\}$, is a list satisfying assignment of $\Pi$. For every constraint $\phi_i=\langle (w_{i,1},w_{i,2}),R_i\rangle$ with $\Gamma_i$ being the image of the rectangular mapping, by the property of $(r\cdot k,|\Gamma_i|)$-set gadget and the fact that $|\mathcal S'|\le r \cdot k$, $\mathcal S'$ must include $C_{j}^{(i)}$ and $\overline{C_j^{(i)}}$ for some $1 \le j \le |\Gamma_i|$, which implies that $\mathcal S'$ includes $S_{w_{i,1},a}$ and $S_{w_{i,2},b}$ for some $(a,b) \in R_i$, and thus $\phi_i$ is list satisfied.

%This is because by the property of $(r\cdot k,|\Gamma_i|)$ and the fact that , the only way to cover $\mathcal M^{(i)}$ is to pick $P^{(i)}_{u,1}$ and $P^{(i)}_{u,2}$ simultaneously for some $u \in C_i$. Thus, we must have included $S_{w_{i,1},a}$ and $S_{w_{i,2},b}$ in $\mathcal S'$ for some $(a,b)\in R_i$, which means the constraint $\phi_i$ is list satisfied. Therefore, if $\Pi$ is not $r$-average-list satisfiable, the minimum number of sets need to cover $U$ is greater than $r\cdot k$.

From \Cref{lem:construct-set-gadget}, a $(r\cdot k,|\Gamma_i|)$-set gadget can be constructed in time $\text{poly}(|\Gamma_i|,2^{r \cdot k})$. The whole reduction runs in FPT time while preserving $k=|X|$. Thus, an $f(k)\cdot n^{O(1)}$ time algorithm for $r$-approximating \kexact{} would give an $f(k)\cdot |\Sigma|^{O(1)}$ algorithm for the $r$-average-list satisfiability of 2CSP with rectangular relations, contradicting the strengthened version of Average Baby PIH.
\end{proof}

We remark that the direct product construction does have rectangular relations, although it does not directly give us Average Baby PIH, in view of Example~\ref{ex:1}.
